\section{Reproducibility and Artifact Details}
\label{app:reproducibility}
The appendix should carry details that help a reader reproduce or audit the
work without interrupting the main argument. Keep assumptions and evidence
needed to assess the central claim in the main paper.

\begin{table}[h]
  \caption{Suggested artifact manifest. Fill in exact values for the study.}
  \centering\small
  \begin{tabularx}{\linewidth}{@{}lX@{}}
    \toprule
    \textbf{Item} & \textbf{What to record} \\
    \midrule
    Source & Commit, license, and commands for each experiment \\
    Environment & OS, accelerator, driver, compiler, and library versions \\
    Data & Provenance, preprocessing, splits, and access conditions \\
    Models & Checkpoints, tokenizer, precision, and inference settings \\
    Compute & Budget, seeds, run count, and failure handling \\
    Evidence & Raw measurements, analysis scripts, and expected outputs \\
    \bottomrule
  \end{tabularx}
\end{table}

\subsection{Example configuration}
\begin{lstlisting}[language=STARPython,caption={An illustrative configuration
schema. Values must be chosen for the actual experiment.},label={lst:config}]
experiment = {
    "model": "replace-with-checkpoint",
    "seed": 0,
    "precision": "float32",
    "split": "held-out-test",
    "record_raw_outputs": True,
}
\end{lstlisting}

\subsection{Editing the manuscript}
Edit \code{paper-info-single.tex} to change the title, authors, affiliation,
keywords, and short running title. Edit the files under
\code{sections/single/} to replace the sample content. Keep new notation in
\code{macros.tex}, references in \code{references.bib}, and illustrations in
\code{figures/}. The shared class controls the visual design.

The bibliography uses BibTeX and author--year citations in this layout.
Use \verb|\citet{key}| for a textual citation and \verb|\citep{key}| for a
parenthetical citation. Use \verb|\autoref{label}| to create typed references
to sections, equations, figures, tables, and algorithms.

\subsection{Modes and publication workflow}
Add \code{anonymous} to the document-class options to hide the lab emblem,
authors, affiliation, email addresses, acknowledgments, and identifying
PDF metadata managed by this class. Inspect the manuscript content,
citations, links, and external assets separately for identity disclosures.
Add \code{monochrome} for a grayscale version of the class palette and logo.

Use the lab style for preprints, technical reports, and internal drafts.
For a specific venue, use its required class and current author instructions.
The section files, bibliography, equations, and vector figures are designed
to be reusable, but class-specific helpers need to be replaced when moving
to another style. The package's README explains this handoff.
